\documentclass[11pt]{article}
\usepackage[utf8]{inputenc}
\usepackage[T2A]{fontenc}
\usepackage[russian]{babel}
\usepackage{amsmath,amssymb}
\usepackage[margin=2.5cm]{geometry}
\usepackage{xcolor}
\definecolor{linkblue}{RGB}{30,80,160}
\usepackage[colorlinks=true, urlcolor=linkblue, linkcolor=black, citecolor=black, breaklinks=true]{hyperref}
\usepackage{booktabs}
\usepackage{enumitem}

\title{Модель нарушителей и модель угроз}
\author{}
\date{22 сентября 2026}

\begin{document}
\maketitle

\section{Роли}

\begin{description}
  \item[Роботы (участники)] $P=\{1,\dots,N\}$, каждый держит приватное значение
    $x_i$ (non-conformity score, дельта весов и т.\,п.).
  \item[Домохозяйства] $H_1,\dots,H_m$ — непересекающиеся группы роботов,
    $P=\bigsqcup_j H_j$ (см.\ документ~06, раздел про compartmented secret
    sharing).
  \item[Сервер(а)] Один сервер $S$ (разделы 1--3 документа~06) или несколько
    $S_1,\dots,S_k$ (Механизм~4, multi-server). $S \cap P = \varnothing$ всегда.
  \item[Внешний наблюдатель] Не участвует в протоколе; видит только выходы
    развёрнутой политики (действия робота) — релевантен основному
    security-роадмапу (membership inference, реконструкция), не MPC-слою напрямую.
  \item[Дилер] Сторона, устанавливающая начальные PRF-ключи / доли на этапе
    настройки протокола (может быть распределённым процессом, не обязательно
    отдельным лицом).
\end{description}

\section{Активы}

\begin{itemize}
  \item \textbf{Защищаем:} индивидуальное значение $x_i$ каждого робота;
    принадлежность робота к домохозяйству (опционально, ниже приоритет).
  \item \textbf{Не защищаем по дизайну:} сам агрегат $\sum_i x_i$ — сервер обязан
    его узнать, это цель протокола, а не утечка.
\end{itemize}

\section{Классы противника}

\begin{enumerate}[label=\textbf{A\arabic*.}, leftmargin=1.6cm]
  \item \textbf{Honest-but-curious единственный сервер.} Строго следует протоколу,
    но пытается извлечь больше, чем положено, из того, что видит. Базовая модель
    для PRSS/Bonawitz (документ~06, Механизм~1).
  \item \textbf{Malicious единственный сервер.} Отклоняется от протокола: например,
    в threshold-HE применяет ключ к \emph{индивидуальному} $c_i$ вместо агрегата,
    или подделывает сообщение об отвалившемся участнике, чтобы получить лишние
    доли. Требует verifiability, а не только маскирования (документ~06, раздел
    \guillemotleft PRSS и Paillier совместно\guillemotright).
  \item \textbf{Сговор $<t$ участников.} По определению порога схемы им не
    полагается узнать ничего — это не угроза, а формальная граница корректности
    схемы.
  \item \textbf{Сговор $\ge t$ независимых участников.} Успешная атака по
    построению схемы — параметр риска, управляемый выбором $t$, а не баг протокола.
  \item \textbf{Домохозяйственная коалиция.} Несколько роботов одного дома
    формально независимы по сырому счёту $|S|$, но фактически — один источник
    доверия (общий Wi-Fi, общий владелец). Наивный $(t,N)$-Shamir этого не
    учитывает; описано через домохозяйственный след $h(S)$ (документ~06,
    Механизм~3).
  \item \textbf{Byzantine-робот.} Не honest-but-curious — активно шлёт испорченное
    $x_i$ (data poisoning) при формально корректном соблюдении MPC-протокола
    (доли/маски валидны, содержимое — нет). Secure aggregation сама по себе это не
    ловит; нужен отдельный слой (FedGT/group testing, документ~06, \guillemotleft
    Направление улучшения\guillemotright).
  \item \textbf{Сговор серверов (multi-server, Механизм~4).} При $m$ серверах
    приватность держится, пока число сговорившихся серверов меньше порога (1 из 2
    у Prio, честное большинство у ABY3-стиля). Сговор всех $m$ серверов ломает
    защиту полностью — фундаментальное ограничение модели, не инженерный дефект.
  \item \textbf{Злонамеренный дилер.} Компрометация этапа установки PRF-ключей/
    долей ломает все последующие раунды. Verifiable-схемы (документ~06) защищают
    именно от этого класса.
\end{enumerate}

\section{Формальное условие нарушения приватности}

\begin{center}
\begin{tabular}{@{}p{5.2cm}p{9.8cm}@{}}
\toprule
Механизм & Условие успешной атаки \\
\midrule
Single-server PRSS/Shamir & $|S_{\mathrm{colluding}}|\ge t$ \\
Single-server + домохозяйства & $h(S_{\mathrm{colluding}})\ge \tau$ (число
  домохозяйств, не сырое число роботов) \\
Threshold-HE без маски & $\ge t$ участников сговариваются \emph{и} применяют ключ
  к индивидуальному $c_i$, а не к агрегату \\
Threshold-HE + PRSS-маска & то же самое, но расшифровка отдельного $c_i$ даёт
  $x_i+z_i$, не $x_i$ — атака не даёт результата без дополнительного взлома
  PRF-структуры \\
Multi-server (Prio/ABY3) & сговор всех $m$ серверов (или превышение порога честного
  большинства) \\
Byzantine-детекция отсутствует & 1 испорченный робот не пойман FedGT-style
  механизмом \\
\bottomrule
\end{tabular}
\end{center}

\section{Обязательства участников (что должно происходить, чтобы гарантия держалась)}

\subsection*{Роботы}
\begin{itemize}
  \item Честно участвовать в первоначальной генерации PRF-ключей/долей на этапе
    настройки — без verifiable-схемы это не проверяется автоматически, только с
    ней (VSS).
  \item Никогда не публиковать свою долю/ключ вне защищённого канала к серверу(ам).
  \item Использовать настоящие значения $x_i$, не сфабрикованные. \textbf{Это не
    гарантируется криптографией самого MPC-слоя} — требует отдельного
    byzantine-детектора (FedGT-style, документ~06).
\end{itemize}

\subsection*{Сервер(а)}
\begin{itemize}
  \item Не пытаться обрабатывать индивидуальные вклады отдельно от агрегата, даже
    если это технически возможно (см.\ A2) — отклонение от протокола, а не
    легальная операция.
  \item В multi-server модели — быть реально независимыми организациями. Несколько
    серверов одного вендора в разных дата-центрах \textbf{не} создают настоящую
    non-collusion (документ~06, Механизм~4).
\end{itemize}

\subsection*{Оператор системы}
\begin{itemize}
  \item Осознанно выбрать параметры: порог $t$, домохозяйственный порог $\tau$,
    число серверов $m$ — исходя из реалистичной оценки, сколько роботов/серверов
    может быть одновременно скомпрометировано, а не произвольно.
  \item Явно решить, нужна ли verifiability (VSS) — то есть учитывается ли класс
    A2/A8 (malicious сервер/дилер) в модели угроз, или ограничиваемся
    honest-but-curious (A1).
\end{itemize}

\section{Явно вне рамок этого документа}

Атаки класса \guillemotleft внешний наблюдатель\guillemotright{} (membership
inference, реконструкция обучающих данных по выходам развёрнутой политики) —
предмет основного security-роадмапа проекта (\emph{Your Robot Remembers Your
Home}), не MPC-слоя. Этот документ формализует только угрозы внутри protocol'а
агрегации.

\end{document}
